\documentclass{article}
\usepackage{hyperref}

\begin{document}
    \begin{center}
        \Large \textbf{Quick Start - Cheat Sheet} \\ [6pt]
        \normalsize
        Autor: Maximilian Jakob Maag B.Sc. \\
        Version: 1.0
    \end{center}
    \section{Intro}
    Dieses Cheat-Sheet liefert Quick-Start-Infos für den Umgang mit Linux. \\\\
    Als erstes muss eine Distribution ausgewählt werden. Zu den bekanntesten Distributionen zählen Ubuntu, Debian, Arch, Linux Mint. \\
    Das ISO-File der Distribution ist auf der Webseite der Distribution herunterzuladen und auf einen USB-Stick zu flashen. \\
    Unter Windows kann Rufus verwendet werden. Linux und MacOS Nutzer verwenden DD. \\
    Das BIOS/UEFI muss auf der Hardware vorbereitet werden. Secure Boot sollte ausgeschaltet werden.
    Alternativ ist die Anleitung der entsprechenden Distribution zur Erstellung der Secure-Boot-Keys zu befolgen.
    \section{Links}
    \begin{itemize}
        \item \href{https://ubuntu.com}{Ubuntu Webseite} Webseite der Distribution Ubuntu.
        \item \href{https://www.debian.org/index.de.html}{Debian Webseite} Webseite der Distribution Debian.
        \item \href{https://rufus.ie/de/}{Rufus Webseite} Startseite Rufus. Tool, um ISO-File auf USB-Stick zu flashen.
    \end{itemize}
\end{document}